In this subsection we set notation for working in the various categories of interest.
%Our goal here is to provide a gentle introduction so that reader interested in working with $C\ta$-modules has good computational control over the category.
%This means digesting the category on the right-hand-side of \Cref{thm:Cta}
In particular, we aim for a practical understanding of the category
\[ \Mod ( \Sp_{C_2, i2}; \uZt) \otimes_{\Z} \mathrm{IndCoh}(\mathcal{M}_{\mathrm{fg}}) \]
that appears on the right-hand-side of \Cref{thm:Cta}. Proceeding from inside out,
we start by fixing notation for the category of $\Z$-modules:

\begin{dfn}
  Let $\Abh$ denote the abelian category of discrete abelian groups.
  Let $\Ab$ denote the category of $\Z$-modules with its standard $t$-structure,
  so that $\Abh$ is the heart of $\Ab$.
\end{dfn}

%From here this subsection now breaks into three parts.
%First, we discuss the relevant categories of equivariant objects.
%Next, we discuss the relevant category of $\MU_*\MU$-comodules, or equivalently sheaves on the moduli stack of formal groups $\Mfg$.
%Finally, we explain how to combine these two pieces.
%Thankfully, the equivariant and comodule aspects are only weakly coupled so we will be able to give simple descriptions.

%% We begin by recalling the definitions of the relevant equivariant categories. These include (discrete) Mackey functors for $C_2$ and modules over $\uZt$ in Mackey functors.
%% Let us begin by taking the time to carefully set up the various players in the story. We start by discussing some categories of Mackey functors and functors between them.

Next, we describe the $C_2$-equivariant analog of $\Ab$:

\begin{dfn} \label{dfn:mackey}
  Let $M(C_2)^{\heartsuit}$ denote the abelian category of discrete Mackey functors for the group $C_2$,
  which may be explicitly described as follows.
  A Mackey functor $B \in M(C_2) ^{\heartsuit}$ consists of the following:
  \begin{itemize}
    \item abelian groups $B(C_2)$ and $B(\ast)$,
    \item an involution $\sigma : B(C_2) \to B(C_2)$,
    \item homomorphisms $r : B(\ast) \to B(C_2)$ and $t : B(C_2) \to B(\ast)$,
     which satisfy the relations $\sigma \circ r = r$, $t \circ \sigma = t$ and $r \circ t = 1 + \sigma$.
  \end{itemize}
  We will somtimes write $[C_2]$ for the endomorphism $t \circ r$ of $B(*)$.
  We let $M(C_2)$ denote the unbounded derived category of $M(C_2)^{\heartsuit}$ equipped with its standard $t$-structure.\footnote{As defined, for example, in \cite[Definition 1.3.5.8]{HA}.}
\end{dfn}

In our setting, $M(C_2)$ arises naturally because it is equivalent to the category of $\piu_0 ^{C_2} \Ss$-modules in $\Sp_{C_2}$ \cite[Theorem 5.10]{SpectrumMackey}.
This description equips $M(C_2)$ with a symmetric monoidal structure compatible with the $t$-structure. The induced symmetric monoidal structure on $M(C_2)^{\heartsuit}$ is known as the box product and admits an explicit description as in \cite{CpBox}.
\todo{Probably there's a better cite for this, but this one has a nice modern proof.}
\todo{rwb: Its always unnerved me that I don't know an explicit description of the monoidal structure on Mackey functors.}

\begin{exm}
  Given a discrete abelian group $A$, we may consider the constant Mackey functor $\underline{A} \in M(C_2)^{\heartsuit}$. This Mackey functor is determined by $\underline{A} (C_2) = \underline{A} (\ast) = A$, $r = \sigma = \mathrm{id}_{A}$ and $t=2 \cdot \mathrm{id}_{A}$.
  This construction provides an exact colimit-preserving functor
  \[\_ : \Abh \to M(C_2)^{\heartsuit}.\]
  Examining the explicit description of the box product in \cite{CpBox}, we find that $\_$ acquires the structure of a (non-unital) tensor-product preserving lax symmetric monoidal functor.
  In particular, we obtain a commutative algebra $\uZ$ in $M(C_2)^{\heartsuit}$.
\end{exm}
\todo{rwb: Given we describe nothing about the monoidal structure, why is this functor lax-mon?}

\begin{dfn}
  Let $\underline{\mathrm{Ab}}$ denote the the symmetric monoidal category of $\uZ$-modules in $M(C_2)$ (or, equivalently, $\Sp_{C_2}$), equipped with its standard $t$-structure. This is equivalent to the derived category of $\Mod(M(C_2)^{\heartsuit}; \uZ)$; in particular, we have $\uAbh = \Mod(M(C_2) ^{\heartsuit}; \uZ)$.
\end{dfn}

\begin{rmk} \label{rmk:underline}
  Since $\Z$ is the unit of $\Abh$, we learn that $\_$ factors through an exact colimit-preserving symmetric monoidal functor% \footnote{A priori this functor is only lax symmetric monoidal, but exact colimit-preserving lax symmetric monoidal functors out of the category of abelian groups which preserve the unit are automatically symmetric monoidal.}
  \[\_ : \Abh \to \uAbh.\]
  Since $\_$ is exact, colimit-preserving and symmetric monoidal it extends uniquely to a colimit-preserving symmetric monoidal functor $\_ : \Ab \to \uAb$ of derived categories. Moreover, the pair of functors $A \to A(*)$ and $A \to A(C_2)$ extend to limit-preserving functors on the level of the derived category. Since this pair of functors is jointly conservative, we can use the fact that $\underline{A}(*) \cong A$ and $\underline{A}(C_2) \cong A$ to conclude that $\_$ preserves limits as well.
 
%% \todo{Is it symmetric monoidal? If so, reason/cite? rwb: Apply BBL to the algebra map Sp --> uAb and you obtain the same symm mon functor in from Ab. Is this worth stating? rwb (later): yes, b/c nobodies written down a lax tensor product of infinity categories we need the functor to have been monoidal.}
\end{rmk}

Before proceeding to explicitly describe $\IndCoh(\Mfg)$, we record a few lemmas for later use. 

\begin{lem} \label{lem:const-mackey}
  The functor $\_ : \Abh \to \uAbh$ is fully faithful and left adjoint to the functor
  $B \mapsto B(\ast)$.
  Its essential image is therefore a coreflective subcategory of $\uAbh$, with coreflector given by the counit map $\underline{B(\ast)} \to B$.
  In particular, for any $B$ in the essential image of $\_$, there are canonical isomorphisms
  \[B \cong \underline{B(\ast)} \cong \underline{B(C_2)}.\]
\end{lem}

\begin{proof} Clear.
\end{proof}

\begin{lem} \label{lem:uabh-image}
  The forgetful functor $\uAbh \to M(C_2)^{\heartsuit}$ is fully faithful, with image spanned by the Mackey functors $B$ for which $[C_2]=2$ as endomorphisms of $B(\ast)$.
\end{lem}

\begin{proof}
  The unit of $M(C_2)^{\heartsuit}$ is the Burnside Mackey functor $\underline{\mathcal{A}}$, whose values are given by $\underline{\mathcal{A}} (C_2) = \Z$ and $\underline{\mathcal{A}} (\ast) = \Z[[C_2]]/([C_2]^2 - 2[C_2])$.
  As a commutative algebra in $M(C_2)^{\heartsuit}$, $\uZ$ is the quotient of $\underline{\mathcal{A}}$ by the relation $[C_2]=2$.
  The result follows.
\end{proof}

The following definition and construction sum up what we need to know about categories of quasicoherent and ind-coherent sheaves on $\Mfg$:

\begin{dfn} \label{dfn:IndCoh}
  We let $\Mfg$ denote the moduli stack of formal groups and let $\QCoh (\Mfg)$ denote the category of quasicoherent sheaves on $\Mfg$.
  It is equivalent to the derived category of evenly graded $\MU_* \MU$-comodules \cite[Remarks 2.38 and 3.14]{GoerssMfg}.
%
  %Let $\QCoh(\Mfg)_{i2}$ denote the category of modules over the $2$-completion of the unit.

  Let $\mathcal{D} \subset \QCoh(\Mfg)$ denote the thick subcategory generated by the sheaves $\omega_{\mathbb{G}/\Mfg} ^{\otimes k}$ for $k \in \Z$, where $\omega_{\mathbb{G} / \Mfg}$ is the sheaf of invariant differentials on the universal formal group $\mathbb{G} / \Mfg$. We define $\IndCoh(\Mfg) \coloneqq \Ind(\mathcal{D})$.
  %It follows from \cite[Remark 4.30]{LocDual} that $\IndCoh(\Mfg)$ is equivalent to .
  The category $\IndCoh(\Mfg)$ is equivalent to Hovey's category of evenly graded stable comodules $\mathrm{Stable}_{\MU_* \MU}$, c.f. \cite[Remark 4.30]{LocDual}.
  %
  %Finally, we let $\IndCoh(\Mfg)_{i2}$ denote the category of modules over the $2$-completion of the unit.

  Both $\QCoh(\Mfg)$ and $\IndCoh(\Mfg)$ are naturally stable presentably symmetric monoidal categories and come equipped with compatible $t$-structures whose hearts are equivalent to the $1$-category of $\MU_* \MU$-comodules. \todo{Come back to this and remark that D = Perf and justify that Indcoh = Indperf here.}
\end{dfn}

\begin{cnstr}
  Since $\QCoh(\Mfg)$ is equivalent to the derived category of a Grothendieck abelian category, it is equipped with the structure of a colimit-preserving symmetric monoidal functor $\Ab \to \QCoh(\Mfg)$.
  Letting $\Ab ^{\mathrm{fin}}$ denote the full subcategory of compact objects, this restricts to an exact symmetric monoidal functor $\Ab ^{\mathrm{fin}} \to \D$, where $\D$ is defined as in \Cref{dfn:IndCoh}.
  Taking $\Ind$ and using the fact that $\Ab$ is compactly generated, we obtain a colimit-preserving symmetric monoidal functor
  \[\Ab \simeq \Ind(\Ab ^{\mathrm{fin}}) \to \Ind(\D) = \IndCoh(\Mfg).\]

  Using this functor and the functor $\Ab \xrightarrow{\_} \uAb \xrightarrow{- \otimes_{\uZ} \uZt} \uAb_{i2}$, we are able to define the stable presentably symmetric monoidal category
  \[\uAb_{i2} \otimes_{\Z} \IndCoh(\Mfg),\]
  which appeared in \Cref{thm:Cta} with slightly different notation.
\end{cnstr}

The core of our understanding of $\uAb_{i2} \otimes_{\Z} \IndCoh(\Mfg)$ rests on the following lemma, where we put a $t$-structure on this category whose heart we can describe explicitly.

\begin{lem} \label{lem:Cta-t-structure}
  There exists a $t$-structure on $\uAb_{i2} \otimes_{\Ab} \IndCoh(\Mfg)$ with the following properties:
  \begin{enumerate}
    \item It is compatible with the symmetric monoidal structure, accessible, compatible with filtered colimits and right complete.
    \item The natural functor
      \[\uAb_{i2, \geq 0} \otimes _{\Ab_{\geq 0}} \IndCoh(\Mfg)_{\geq 0} \to \uAb_{i2} \otimes_{\Ab} \IndCoh(\Mfg)\]
      is fully faithful with essential image the connective objects.
    \item The natural functor
      \[\uAb_{i2} ^{\heartsuit} \otimes_{\Ab^{\heartsuit}} \IndCoh(\Mfg)^{\heartsuit} \to \uAb_{i2} \otimes_{\Ab} \IndCoh(\Mfg)\]
      is fully faithful with essential image the heart of $\uAb_{i2} \otimes_{\Ab} \IndCoh(\Mfg)$.
  \end{enumerate}
\end{lem}

Using the explicit description of $\uAbh$ given by \Cref{dfn:mackey} and \Cref{lem:uabh-image} and the fact that $\IndCoh(\Mfg)^{\heartsuit}$ is equivalent to the category of evenly-graded $\MU_* \MU$-comodules, we obtain the following description of
\[\uAb_{i2}^{\heartsuit} \otimes_{\Ab^{\heartsuit}} \IndCoh(\Mfg)^{\heartsuit}:\]
this is the category of pairs of $(\MU_2)_*\MU_2$-comodules $A(C_2)$ and $A(\ast)$ equipped with following structures (all of which are compatible with the comodule structure):
\begin{itemize}
  \item maps $r : A(\ast) \to A(C_2)$ and $t : A(C_2) \to A(\ast)$, along with an involution $\sigma : A(C_2) \to A(C_2)$,
  \item which satisfy the relations $\sigma \circ r = r$, $t \circ \sigma = t$, $t \circ r = 2$ and $r \circ t = 1 + \sigma$.
\end{itemize}

In other words, $\uAb^{\heartsuit}_{i2} \otimes_{\Ab^{\heartsuit}} \IndCoh(\Mfg)^{\heartsuit}$ is the category of evenly graded comodules over the Hopf algebroid $(\underline{(\MU_2)_{*}}, \underline{(\MU_2)_* \MU_2})$ in graded $C_2$-Mackey functors.

\begin{rmk}
  Using the fact that $$((\underline{\MU_{\R,2}})_{*\rho}, (\underline{\MU_{\R,2}})_{*\rho} \MU_{\R,2}) \cong (\underline{(\MU_2)_{2*}}, \underline{(\MU_2)_{2*} \MU_2}),$$ see \cite[Theorems 2.25 and 2.28]{HuKriz}, this category may equally well be described as the category of graded comodules over the Hopf algebroid $((\underline{\MU_{\R,2}})_{*\rho}, \underline({\MU_{\R,2}})_{*\rho} \MU_{\R,2})$ in graded $C_2$-Mackey functors.
\end{rmk}

The proof of \Cref{lem:Cta-t-structure} is not difficult, but it will rely on material from \cite[Appendix C]{SAG} which we presently recount. We begin with a simple method for producing $t$-structures. 

\begin{cnstr}
  Given a pointed presentable category $\CC$, we let 
  $\Sp \otimes \CC \simeq \Sp(\CC)$
  %% = \varprojlim \left( \dots \xrightarrow{\Omega} \CC \xrightarrow{\Omega} \CC \xrightarrow{\Omega} \CC \right)\]
  denote the category of spectrum objects in $\CC$. It is a stable presentable category.

  The category $\Sp \otimes \CC$ admits a $t$-structure with $(\Sp \otimes \CC)_{\geq 0}$ equal to the essential image of the functor $\Sigma^{\infty} : \CC \to \Sp \otimes \CC$.
  This $t$-structure is accessible and right complete.

  Note that if $\CC$ is endowed with the structure of a presentably symmetric monoidal category, then $\Sp \otimes \CC$ naturally inherits this structure and the $t$-structure defined above is compatible with it.
\end{cnstr}

In fact, the following lemma, which follows from \cite[Remark C.3.1.5]{SAG}, demonstrates that this construction is the universal way to produce a (well-behaved) $t$-structure on a presentable stable category. 

\begin{lem} \label{lem:conn-enough}
  Suppose that $(\CC, \CC_{\geq 0})$ is a stable presentably symmetric monoidal category with compatible $t$-structure. Suppose that the $t$-structure is accessible, compatible with filtered colimits and right complete.

  Then there is a natural equivalence of symmetric monoidal categories with compatible $t$-structure
  \[(\CC, \CC_{\geq 0}) \simeq (\Sp \otimes \CC_{\geq 0}, (\Sp \otimes \CC_{\geq 0})_{\geq 0}).\]
\end{lem}

%% \begin{cor} \label{cor:is-sp}
%%   There are natural equivalences of symmetric monoidal categories with compatible $t$-structures:
%%   \[(\Ab,\Ab_{\geq 0}) \simeq (\Sp(\Ab_{\geq 0}), \Sp(\Ab_{\geq 0})_{\geq 0}),\]
%%   \[(\uAb,\uAb_{\geq 0}) \simeq (\Sp(\uAb_{\geq 0}), \Sp(\uAb_{\geq 0})_{\geq 0}),\]
%%   \[(\IndCoh(\Mfg), \IndCoh(\Mfg)_{\geq 0}) \simeq (\Sp(\IndCoh(\Mfg)_{\geq 0}), \Sp(\IndCoh(\Mfg)_{\geq 0})_{\geq 0}).\]
%% \end{cor}

In conclusion, we may as well work with the categories of connective objects. For $\Ab, \uAb_{i2}$, and $\IndCoh(\Mfg)$, these are all Grothendieck prestable categories in the sense of \cite[Definition C.1.4.2]{SAG}.

\begin{rec} \label{rec:groth-ab}
  We let $\Groth_{\infty}$ denote the category of Grothendieck prestable categories \cite[Definition C.3.0.5]{SAG}. By \cite[Theorem C.4.2.1]{SAG}, $\Groth_{\infty}$ is a symmetric monoidal category with tensor product given by the usual tensor product of presentable categories and unit $\Sp_{\geq 0}$.

  Moreover, let $\Groth_{1}$ denote the category of Grothendieck abelian $1$-categories. It is a symmetric monoidal category with tensor product given by the usual tensor product of presentable categories and unit $\Abh$ \cite[Corollary C.5.4.19]{SAG}. There is a functor $\tau_{\leq 0} : \Groth_{\infty} \to \Groth_{1}$ sending a Grothendieck prestable category to its subcategory of discrete objects.
  By \cite[Remark C.5.4.20]{SAG}, the functor $\tau_{\leq 0}$ is symmetric monoidal.

  We summarize this in the following span of symmetric monoidal categories 
  \begin{center}
    \begin{tikzcd}
      & \Groth_{\infty} \ar[dl, "\Sp \otimes -"] \ar[dr, "\tau_{\leq 0}"] \\
      \mathrm{Pr}^{L,\mathrm{Stab}} & & \Groth_{1}.
    \end{tikzcd}
  \end{center}
\end{rec}

\begin{proof}[Proof (of \Cref{lem:Cta-t-structure}).]
  Using \Cref{lem:conn-enough} and \cite[Remark C.4.2.3]{SAG}, we have
  \[\Sp \otimes \left(\uAb_{i2, \geq 0} \otimes_{\Ab_{\geq 0}} \IndCoh(\Mfg)_{\geq 0} \right) \simeq \uAb_{i2} \otimes_{\Ab} \IndCoh(\Mfg).\]
  Moreover, since $\uAb_{i2, \geq 0} \otimes_{\Ab_{\geq 0}} \IndCoh(\Mfg)_{\geq 0}$ is prestable, the $\Sigma^{\infty}$ functor identifies it with $\left(\Sp \otimes \left(\uAb_{i2, \geq 0} \otimes_{\Ab_{\geq 0}} \IndCoh(\Mfg)_{\geq 0} \right)\right)_{\geq 0}$.

  This immediately implies the first two properties, with the exception of compatibility with filtered colimits. This follows from the fact that $\uAb_{i2, \geq 0} \otimes_{\Ab_{\geq 0}} \IndCoh(\Mfg)_{\geq 0}$ is compactly generated, which holds becuase each one of $\uAb_{i2, \geq 0}$, $\Ab_{\geq 0}$ and $\IndCoh(\Mfg)_{\geq 0}$ is compactly generated \cite[Corollary C.6.2.3]{SAG}.

  The third property follows from the fact that $\CC^{\heartsuit} = \tau_{\leq 0} \left( \CC_{\geq 0} \right)$ and the fact that $\tau_{\leq 0}: \Groth_{\infty} \to \Groth_1$ is symmetric monoidal.
\end{proof}

%\begin{dfn}
%  Let $\CC$ and $\D$ denote stable presentably symmetric monoidal $\infty$-categories equipped with $t$-strutures generated by connective objects $\CC_{\geq 0}$ and $\D_{\geq 0}$. We make the following additional assumptions on the $t$-structure:
%  \begin{itemize}
%    \item We have $\o_\CC \in \CC_{\geq 0}$ and $\o_{\D} \in \D_{\geq 0}$.
%    \item The full subcategories $\CC_{\geq 0}$ and $\D_{\geq 0}$ are closed under the tensor product.
%    \item The full subcategories $\CC_{\geq 0}$ and $\D_{\geq 0}$ of coconnective objects are closed under filtered colimits.
%  \end{itemize}
%\end{dfn}
%
%Remark C.5.4.20
%
%Finally, we define a $t$-structure on $\uAb \otimes_\Z \IndCoh(\Mfg)$ by setting $\left( \uAb \otimes_\Z \IndCoh(\Mfg) \right) _{\geq 0}$ to be the category generated under colimits and extensions by the image of the restriction of the functor $m : \uAb \times \IndCoh(\Mfg) \to \uAb \otimes_{\Z} \IndCoh(\Mfg)$ to $\uAb_{\geq 0} \times \IndCoh(\Mfg)_{\geq 0}$.
%
%
%We may now make sense of the tensor product on the right-hand side of the isomorphism in \Cref{thm:Cta}. By the description of $\uAb$ in \Cref{prop:uA-pre}, we see that there is an equivalence
%\[\Mod(\Sp_{C_2}; \uZ) \otimes_{\Z} \IndCoh(\Mfg) = \uAb \otimes_{\Z} \IndCoh(\Mfg) \simeq \Fun^{\times} (\uA, \IndCoh(\Mfg)).\]
%



%
%The usual $t$-structure on $\IndCoh(\Mfg)$ induces a $t$-structure on $\Fun^{\times} (\uA, \IndCoh(\Mfg))$ with heart $\Fun^{\times} (\uA, \IndCoh(\Mfg)^{\heartsuit})$. Unpacking the definitions, we find that an object of this heart consists \todo{up to p-completion nonsense...} of a pair of $\MU_*\MU$-comodules $A(C_2)$ and $A(\ast)$ equipped with following strucutre (all of which is compatible with the comodule structure),
%\begin{itemize}
  %\item an involution, $\sigma$, on $A(C_2)$,
  %\item restriction and transfer maps, $r: A(\ast) \to A(C_2)$ and $t: A(C_2) \to A(\ast)$,
%\end{itemize}
%subject to compatabilities $r \sigma = r$, $\sigma t = t$, $rt = 2$ and $tr = 1 + \sigma$.
%
%
%
%Now we begin our task of understanding the tensor product on the right-hand-side of the isomorphism in the statement of [thm]. Since $\uAb$ can be described as product preserving presheaves of $\Z$-modules on $\uA$ we can correspondingly describe $\Mod ( \Sp_{C_2}; \uZ) \otimes_{\Z} \mathrm{IndCoh}(\mathcal{M}_{\mathrm{fg}})$ as product preserving presheaves on $\uA$ with values in $\mathrm{IndCoh}(\mathcal{M}_{\mathrm{fg}})$.  
%
%By declaring that external tensors of connective objects are connective we obtain a $t$-structure on $\Mod ( \Sp_{C_2}; \uZ) \otimes_{\Z} \mathrm{IndCoh}(\mathcal{M}_{\mathrm{fg}})$. The sheaf description gives us a simple description of the heart of this $t$-structure. 
%An object of $(\uAb \otimes_{\Ab} \mathrm{IndCoh}(\mathcal{M}_{\mathrm{fg}}))^\heartsuit$
%consists of a pair of $\MU_*\MU$-comodules $A_e$ and $A_{C_2}$ equipped with following strucutre (all of which is compatible with the comodule structure),
%\begin{itemize}
%\item an involution, $\sigma$, on $A_e$,
%\item restriction and transfer maps, $r: A_{C_2} \to A_e$ and $t: A_e \to A_{C_2}$,
%\end{itemize}
%subject to compatabilities $r \sigma = r$, $\sigma t = t$, $rt = 2$ and $tr = 1 + \sigma$.
%
%Using this description we can conclude that the statement about the trigraded homotopy ring in [thm] follows from equivalence of categories given earlier in [thm]. \todo{Actually prove this and come back to it later.}
%We will prove the remainder of [thm] over the course of the next two subsections.
%In [subsec], we will use general nonsense to produce a comparison functor between the two categories in question and reduce [thm] to a single, more technical lemma.
%In [subsec], we will use the special properties of $t$-structures to prove this lemma.
%\todo{Update this at the end to take into account the fact that there are two theorems now.}
%
%% \begin{rmk}\todo{Move this elsewhere}
%%   Although we have mentioned Mackey functors in connection with >>> it is not true that >>> is the category of $\MU_*\MU$-comodule valued Mackey functors. We suspect that if one studied $C_2$ equivariant motives, i.e. $\SH(B\mu_2)$, that a corresponding story would play out with $\MU_\R$ replaced by $\MU_{C_2}$. In this case the the modules over associated [tau] element would actually be comodule valued Mackey functors. 
%% \end{rmk}
%
%The proof of this theorem will occupy the next several pages.
%We begin with several lemmas whose proofs are short and abstract (Lemmas \ref{lem:gr-of-comparison}, \ref{lem:MU-koszul} and \ref{lem:uMU-koszul}).
%These lemmas puts us in a position where \Cref{lem:gr-R-bullet-computation} finishes the proof.
%The proof of \Cref{lem:gr-R-bullet-computation} is much more concrete, it proceeds via a direct examination of the object $R_\bullet$ constructed in [location].
